\chapter{热力学势、配分函数与响应的核心推导}

\section{基本关系与 Legendre 变换}
简单可压缩单组分系统的能量表象取
$U=U(S,V,N)$：
\[
\dd U=\left(\pdv US\right)_{V,N}\dd S
+\left(\pdv UV\right)_{S,N}\dd V
+\left(\pdv UN\right)_{S,V}\dd N.
\]
用温度、压强、化学势的热力学定义：
\[
T=\left(\pdv US\right)_{V,N},\quad
P=-\left(\pdv UV\right)_{S,N},\quad
\mu=\left(\pdv UN\right)_{S,V},
\]
得到
\[
\boxed{\dd U=T\dd S-P\dd V+\mu\dd N}.
\]
实验常控制 $T,P,\mu$ 而非 $S,V,N$，故作 Legendre 变换：
\[
F=U-TS,\quad
\dd F=-S\dd T-P\dd V+\mu\dd N,
\]
\[
G=U-TS+PV,\quad
\dd G=-S\dd T+V\dd P+\mu\dd N,
\]
\[
\Omega=U-TS-\mu N,\quad
\dd\Omega=-S\dd T-P\dd V-N\dd\mu.
\]
每个势的自然变量直接由微分读出。

\section{Maxwell 关系逐一得到}
例如 $F(T,V,N)$：
\[
S=-\left(\pdv FT\right)_{V,N},\qquad
P=-\left(\pdv FV\right)_{T,N}.
\]
混合偏导可交换：
\[
\pdv{}{V}\left(\pdv FT\right)
=\pdv{}{T}\left(\pdv FV\right),
\]
所以
\[
\boxed{\left(\pdv SV\right)_T
=\left(\pdv PT\right)_V}.
\]
同理由 $G(T,P,N)$ 得
\[
\boxed{\left(\pdv SP\right)_T
=-\left(\pdv VT\right)_P}.
\]
由 $U(S,V,N)$ 和焓 $H=U+PV$ 还可得到另外两组。记忆方法不是背符号，而是写出势的微分再交叉求导。

\section{稳定性条件}
固定 $T,V,N$ 时平衡使 Helmholtz 自由能取极小。对温度涨落，
\[
C_V=T\left(\pdv ST\right)_V
=-T\left(\frac{\partial^2 F}{\partial T^2}\right)_V\ge0.
\]
固定温度的机械稳定性要求压缩体积时压强上升：
\[
\left(\pdv PV\right)_T<0,\qquad
\kappa_T=-\frac1V\left(\pdv VP\right)_T>0.
\]
这些条件等价于适当热力学势 Hessian 的凸性/凹性。

\section{热库展开推导正则分布}
系统微态 $i$ 能量 $E_i$，总系统孤立：
\[
P_i\propto\Omega_R(E_{\rm tot}-E_i).
\]
热库巨大，展开熵
$S_R=\kb\ln\Omega_R$：
\[
S_R(E_{\rm tot}-E_i)
=S_R(E_{\rm tot})
-\left(\pdv{S_R}{E_R}\right)E_i
+O(E_i^2/E_R).
\]
$\partial S_R/\partial E_R=1/T$：
\[
\Omega_R(E_{\rm tot}-E_i)
\propto e^{-E_i/\kb T}.
\]
归一化：
\[
\boxed{P_i=\frac{e^{-\beta E_i}}Z},\qquad
Z=\sum_ie^{-\beta E_i}.
\]
二阶项被忽略要求系统远小于热库；保留它可描述有限热库修正。

\section{最大熵法}
最大化
\[
\Phi=-\kb\sum_ip_i\ln p_i
-\alpha(\sum_ip_i-1)
-\gamma(\sum_ip_iE_i-U).
\]
对每个 $p_i$：
\[
\pdv\Phi{p_i}
=-\kb(\ln p_i+1)-\alpha-\gamma E_i=0,
\]
\[
p_i=C e^{-(\gamma/\kb)E_i}.
\]
用热力学关系识别 $\gamma/\kb=\beta$，归一化得同一正则分布。该推导表明：给定平均能量时，正则分布是熵最大的分布。

\section{配分函数生成热力学量}
\[
\pdv{\ln Z}{\beta}
=\frac1Z\sum_i(-E_i)e^{-\beta E_i}
=-\langle E\rangle,
\]
\[
\boxed{U=-\partial_\beta\ln Z}.
\]
熵
\[
S=-\kb\sum_ip_i\ln p_i
=\kb\ln Z+\kb\beta U.
\]
故
\[
F=U-TS=-\kb T\ln Z.
\]
二阶导数：
\[
\partial_\beta^2\ln Z
=\langle E^2\rangle-\langle E\rangle^2.
\]
又
$\partial_\beta=-\kb T^2\partial_T$：
\[
\boxed{\langle(\Delta E)^2\rangle=\kb T^2C_V}.
\]

\section{经典理想气体配分函数}
单粒子平动配分函数
\[
Z_1=\frac1{h^3}\int_V\dd^3q
\int\dd^3p\,e^{-\beta p^2/2m}.
\]
三个 Gaussian 积分
\[
\int_{-\infty}^\infty e^{-\beta p_x^2/2m}\dd p_x
=\sqrt{\frac{2\pi m}{\beta}},
\]
所以
\[
Z_1=\frac V{h^3}(2\pi m\kb T)^{3/2}
=\frac V{\lambda_T^3}.
\]
全同经典粒子消除置换重复计数：
\[
Z_N=\frac{Z_1^N}{N!}.
\]
于是
\[
P=-\left(\pdv FV\right)_{T,N}
=\kb T\pdv{\ln Z_N}{V}
=\frac{N\kb T}{V}.
\]
得到 $PV=N\kb T$。Stirling 公式还可得到 Sackur--Tetrode 熵。

\section{Bose 与 Fermi 分布}
巨配分函数按单粒子态分解：
\[
\Xi=\prod_i\Xi_i,\qquad
\Xi_i=\sum_{n_i}e^{-\beta(\epsilon_i-\mu)n_i}.
\]
玻色子 $n_i=0,1,\ldots$：
\[
\Xi_i^{B}=\frac1{1-e^{-\beta(\epsilon_i-\mu)}}.
\]
费米子 $n_i=0,1$：
\[
\Xi_i^{F}=1+e^{-\beta(\epsilon_i-\mu)}.
\]
平均占据可由
\[
\bar n_i=\frac1\beta\pdv{\ln\Xi_i}{\mu}
\]
求得：
\[
\boxed{\bar n_i^B=
\frac1{e^{\beta(\epsilon_i-\mu)}-1}},
\quad
\boxed{\bar n_i^F=
\frac1{e^{\beta(\epsilon_i-\mu)}+1}}.
\]

\section{Sommerfeld 展开的关键换元}
考虑
\[
I=\int_0^\infty\phi(E)f(E)\dd E.
\]
分离零温阶跃：
\[
I=\int_0^\mu\phi(E)\dd E+
\int_0^\infty\phi(E)[f(E)-\Theta(\mu-E)]\dd E.
\]
第二项局域在 $|E-\mu|\sim\kb T$。令
$x=(E-\mu)/\kb T$，Taylor 展开
\[
\phi(\mu+\kb Tx)=\phi(\mu)+
\kb Tx\phi'(\mu)+\cdots.
\]
修正核为奇函数，常数项积分为零；首个非零矩
\[
\int_{-\infty}^\infty x
[f(x)-\Theta(-x)]\dd x=\frac{\pi^2}{6}.
\]
故
\[
\boxed{I=\int_0^\mu\phi(E)\dd E+
\frac{\pi^2}{6}(\kb T)^2\phi'(\mu)+O(T^4)}.
\]
对粒子数和能量分别使用，得到
$C_V=(\pi^2/3)\kb^2g(E_F)T$。

\section{Landau--Ginzburg 与关联长度}
\[
F[m]=\int\dd^dr\left[
\frac r2m^2+\frac u4m^4+\frac\kappa2(\nabla m)^2-hm
\right],
\quad r=a(T-T_c).
\]
泛函变分：
\[
\delta F=\int\dd^dr[
rm+um^3-h-\kappa\nabla^2m]\delta m
\]
（梯度项分部积分并令边界变分消失）。平衡方程
\[
\boxed{-\kappa\nabla^2m+rm+um^3=h}.
\]
$T>T_c$、小涨落时忽略 $m^3$。Fourier 变换：
\[
(\kappa q^2+r)m_{\vb q}=h_{\vb q},
\qquad
\chi(q)=\frac1{r+\kappa q^2}.
\]
写成
\[
\chi(q)=\frac{\chi(0)}{1+q^2\xi^2},
\qquad
\boxed{\xi=\sqrt{\kappa/r}}.
\]
$T\to T_c$ 时关联长度发散，解释临界涨落为何跨越多尺度。

\section{Kubo 公式从相互作用绘景推出}
$H(t)=H_0-f(t)B$。相互作用绘景密度矩阵满足
\[
i\hbar\dot\rho_I=[H_I'(t),\rho_I].
\]
一阶积分：
\[
\rho_I(t)=\rho_0-\frac{i}{\hbar}
\int_{-\infty}^t[H_I'(t'),\rho_0]\dd t'.
\]
观测量变化
\[
\delta\langle A(t)\rangle
=\mathrm{Tr}[(\rho_I(t)-\rho_0)A_I(t)].
\]
代入 $H'=-fB$，利用迹循环性：
\[
\delta\langle A(t)\rangle
=\frac{i}{\hbar}\int_{-\infty}^t
f(t')\langle[B_I(t'),A_I(t)]\rangle_0\dd t'.
\]
因此
\[
\boxed{\chi_{AB}(t-t')
=\frac{i}{\hbar}\theta(t-t')
\langle[B(t'),A(t)]\rangle_0}.
\]
阶跃函数是因果性；对易子是量子响应的核心。

\begin{selftestbox}
\begin{enumerate}
\item 为什么 $U$ 对 $V$ 的偏导带负号？
\item 如何不背诵地重建 Maxwell 关系？
\item 热库展开忽略二阶项的条件是什么？
\item 最大熵法的两个乘子分别约束什么？
\item $F=-\kb T\ln Z$ 如何推出？
\item Gibbs 因子为何是 $1/N!$？
\item 玻色化学势为何不能高于基态能量？
\item Sommerfeld 展开为何没有线性温度项？
\item 关联长度怎样从 $\chi(q)$ 读出？
\item Kubo 公式中对易子的次序如何确定？
\end{enumerate}
\end{selftestbox}
